% Part: computability
% Chapter: computability-theory
% Section: def-functions-self-reference

\documentclass[../../../include/open-logic-section]{subfiles}

\begin{document}

% SEGMENT OLP-0251-S01

\olfileid{cmp}{thy}{slf}
\olsection{தன்குறிப்பைப் பயன்படுத்திச் சார்புகளை வரையறுத்தல்}

சார்புகளை அவற்றைக்கொண்டே வரையறுக்க முடிவது பொதுவாகப் பயனுள்ளது.
எடுத்துக்காட்டாக, கணிக்கத்தக்க சார்புகள் $k$, $l$,~$m$
கொடுக்கப்பட்டால், ஒவ்வொரு~$y$ க்கும் பின்வரும் சமன்பாட்டை நிறைவேற்றும்
பகுதி கணிக்கத்தக்க சார்பு~$f$ ஒன்று உள்ளது என்று நிலைப்புள்ளித்
துணைத்தேற்றம் கூறுகிறது:
\[
f(y) \simeq
\begin{cases}
k(y) & \text{$l(y) = 0$ எனில்} \\
f(m(y)) & \text{இல்லையெனில்.}
\end{cases}
\]
இன்னும் குறிப்பாக, $f$ பின்வருமாறு பெறப்படுகிறது:
\[
g(x,y) \simeq
\begin{cases}
k(y) & \text{$l(y) = 0$ எனில்} \\
\cfind{x}(m(y)) & \text{இல்லையெனில்}
\end{cases}
\]
என்ற சார்பை எடுத்துக்கொண்டு, $\cfind{e}(y) = g(e,y)$ என அமையும்
சுட்டெண்~$e$ ஐக் கண்டுபிடிக்க நிலைப்புள்ளித் துணைத்தேற்றத்தைப்
பயன்படுத்துகிறோம்.

ஒரு திட்டவட்டமான எடுத்துக்காட்டாக, ``மீப்பெருப் பொதுவகுத்தி''
சார்பு $\fn{gcd}(u,v)$ ஐப் பின்வருமாறு வரையறுக்கலாம்:
\[
\fn{gcd}(u,v) \simeq
\begin{cases}
v & \text{$u = 0$ எனில்} \\
\fn{gcd}(\fn{mod}(v, u), u) & \text{இல்லையெனில்}
\end{cases}
\]
இங்கு $\fn{mod}(v, u)$ என்பது $v$ ஐ~$u$ ஆல் வகுக்கும்போது
கிடைக்கும் மீதியைக் குறிக்கிறது. நிலைப்புள்ளித் துணைத்தேற்றத்தைப்
பயன்படுத்தினால் $\fn{gcd}$ பகுதி கணிக்கத்தக்கது என்பது கிடைக்கிறது.
(உண்மையில், $y$ ஆனது சோடி $\tuple{u, v}$ ஐக் குறிமுறையாக்குவதாக
எடுத்துக்கொண்டு, இதை மேலுள்ள வடிவத்தில் அமைக்கலாம்.) அதன் பின்~$u$
மீதான தொகுத்தறிதல், $\fn{gcd}$ உண்மையில் முழுமையானது என்று காட்டுகிறது.

மேலே விவாதித்த எடுத்துக்காட்டுகளைவிட மிகவும் நுட்பமான தன்குறிப்பு
வரையறைகளை உருவாக்க முடியும் என்பது உண்மையே. பெரும்பாலான நிரலாக்க
மொழிகள் ஏதோ ஒரு முறையில் சார்புகளை அவற்றைக்கொண்டே வரையறுப்பதை
ஆதரிக்கின்றன. ஒரு சார்பைக் கணிக்கும் படிமுறைக்கான ஓர்
\emph{சுட்டெண்ணைக்} கொண்டு அச்சார்பை வரையறுக்க முடிவதைக் காட்டிலும்
இது சற்றுக் குறைவான வியப்புடையது என்பதைக் கவனிக்க; முழுப்
பொதுத்தன்மையில் பின்னதையே செய்ய நிலைப்புள்ளித் தேற்றம் அனுமதிக்கிறது.

\end{document}
